\chapter{Нерадиальная устойчивость эффективного тетраэдрического вихря}

\section{Постановка вопроса}

Предыдущая глава установила положительность радиального гессиана трёх
профилей $K,a,b$. Этого недостаточно для устойчивости прямой вихревой
нити: отрицательная мода может зависеть от полярного угла или продольной
координаты и потому не пересекать осесимметричное подпространство.

Литературная граница здесь существенна. Конечные вихри нарушения
$SO(3)\to A_4$ с голономиями порядка три построены в
\path{arXiv:2607.12366}. Полный анализ скрученных вихрей в
\path{arXiv:0712.3589} показывает, что длинноволновая нерадиальная мода
может быть отрицательной даже при отсутствии однородной неустойчивости.
Фоновая калибровка и отделение физических мод от калибровочных подробно
разобраны для вихрей в \path{arXiv:1602.09084}; самосопряжённый
спектральный оператор вихря в ортогональной калибровке исследован в
\path{arXiv:2608.10608}.

\section{Эффективная поперечная подсистема}

Внутри уже выведенного анзаца поле $a$ является модулем заряженной
комплексной амплитуды
\begin{equation}
 \phi_0(r,\theta)=a(r)e^{i\theta},
 \qquad
 A_{0\theta}=\frac{K(r)}{r},
 \qquad b_0=b(r).
\end{equation}
Её двумерная энергия имеет вид
\begin{equation}
 E_\perp=\int_{\mathbb R^2}
 \left[
  \frac A2|D_i\phi|^2+\frac B2|\partial_i b|^2
  +\frac G2F_{12}^2+V(|\phi|,b)
 \right]d^2x,
 \label{eq:v6-angular-effective-energy}
\end{equation}
где
\begin{equation}
 A=\frac{128}{81},\qquad
 B=\frac{160}{81},\qquad
 G=\frac{2}{27},
\end{equation}
а многочлен $V$ не вводится заново, а берётся из главы о конечном
профиле.

Для возмущений $\eta=\delta\phi$ и $\alpha_i=\delta A_i$ добавлена
фоновая калибровка
\begin{equation}
 \mathcal G(\eta,\alpha)
 =\partial_i\alpha_i+\frac AG(J\phi_0)\mathbin{\cdot}\eta,
 \qquad
 E_{\mathrm{fix}}=\frac G2\int \mathcal G^2\,d^2x.
 \label{eq:v6-angular-background-gauge}
\end{equation}
Она устраняет чистые калибровочные направления, не выбирая сингулярную
фазу в ядре вихря. Проверяемый оператор действует на пять вещественных
компонент
\begin{equation}
 (\Re\eta,\Im\eta,\alpha_x,\alpha_y,\delta b).
\end{equation}

\section{Декартов спектральный аудит}

Оператор второй вариации вычислен прямым дифференцированием дискретной
энергии \eqref{eq:v6-angular-effective-energy} вместе с
\eqref{eq:v6-angular-background-gauge}. Использованы квадратные области
радиуса $10$ и сетки $24$, $32$, $40$, $48$. На последней сетке
размерность оператора равна $10580$, а среднеквадратичный остаток
дискретного уравнения фона равен
\begin{equation}
 4.3252949\cdot10^{-3}.
\end{equation}
Десять нижних собственных значений равны
\begin{equation}
\begin{split}
 &(0.75626863,\ 0.82077050,\ 1.50107385,\ 1.58734966,\
 1.61617826,\\
 &\hspace{1.5cm}1.65323501,\ 1.69354588,\ 1.71651647,\
 1.75128367,\ 1.76502293).
\end{split}
\end{equation}
Отрицательных мод среди проверенных нет. Первые две моды смягчаются при
измельчении сетки и являются кандидатами на решёточно поднятые переносы,
но их строгое отождествление ещё не получено. Зато третье значение
стабилизируется:
\begin{equation}
 |\lambda_3^{(48)}-\lambda_3^{(40)}|
 =3.5780064\cdot10^{-4}.
\end{equation}

Для прямой однородной нити продольный импульс в этой подсистеме добавляет
$k_z^2$. Поэтому проверенный жёсткий канал имеет дисперсию
\begin{equation}
 \omega^2(k_z)=1.50107385\ldots+k_z^2
\end{equation}
и не создаёт длинноволновой отрицательной моды.

\section{Строгая граница результата}

Полученная положительность относится только к эффективной подсистеме,
которая уже присутствует в профиле $Q+T+B$. Она не включает все пять
компонент поля спина два, все семь компонент поля спина три и все
неабелевы компоненты семейной связности. Кроме того, две мягкие моды пока
не сопоставлены аналитически с переносными касательными.

\begin{nogocriterion}
Нельзя объявлять вихрь полностью устойчивым или отождествлять его с
частицей, пока полный калибровочно фиксированный гессиан всех компонент
$Q$, $T$ и $B$ не имеет отрицательных физических мод и воспроизводит
ровно две поперечные переносные нулевые моды.
\end{nogocriterion}

Нулевая гипотеза следующего гейта состоит в том, что одна из исключённых
неабелевых или тензорных компонент содержит отрицательный канал. Если он
обнаружится, прямой вихрь не является устойчивым материальным носителем.

\begin{protocol}
\begin{enumerate}
 \item нерадиальный оператор эффективной подсистемы: \textbf{построен};
 \item отрицательные моды среди десяти нижних: \textbf{не найдены};
 \item жёсткий спектральный зазор: \textbf{$1.50107385\ldots$};
 \item переносная мягкая пара: \textbf{не отождествлена строго};
 \item полный неабелев гессиан: \textbf{не проверен};
 \item рождение материи: \textbf{не замкнуто};
 \item следующий гейт: \textbf{полный гессиан вихря $Q+T+B$}.
\end{enumerate}
\end{protocol}

Машинный сертификат сохранён в
\path{s2t_v6_bosonic_defect_q_tetrahedral_vortex_angular_stability_gate_results.json}.